\section{Capítulo~\ref{sec:dofa}. Aprendizado com \nt{DOPA}}

Os mecanismos da sinapse (porções, detector de coincidência, cálcio, janelas de plasticidade) e o comportamento da dopamina como erro de predição foram medidos diretamente; que as regras levam ao gradiente e quanto custa a difusão são teoremas; o resultado do aprendizado de escolha é um cálculo com o modelo do livro; o circuito que calcula o erro é uma hipótese.

{\footnotesize
\setlength{\tabcolsep}{3pt}\renewcommand{\arraystretch}{1.15}
\begin{longtable}{>{\raggedright}p{0.5cm}>{\raggedright}p{4.0cm}>{\raggedright}p{5.6cm}>{\raggedright}p{2.3cm}>{\raggedright\arraybackslash}p{1.7cm}}
\toprule
N.º & Proposição & Experimento e o que foi medido & Fonte & Status \\
\midrule\endfirsthead
\toprule
N.º & Proposição & Experimento e o que foi medido & Fonte & Status \\
\midrule\endhead
1 & A retropropagação do erro, na forma que tem nas redes artificiais, não foi encontrada no cérebro & Não há experimento direto: é uma afirmação de ausência. As revisões analisam o que o algoritmo exige (conexões de retorno que repetem as diretas, fase separada) e que aproximações biológicas foram propostas. & \cite{Lillicrap2020, WhittingtonBogacz2019} & hipótese \\
2 & O peso se decompõe em número de locais de liberação, probabilidade de liberação e resposta a uma porção: $w=N_s\,p\,A_1$~\eqref{eq:quantal} & Junção neuromuscular de rã com liberação reduzida: a resposta do destinatário varia em degraus múltiplos da resposta miniatura espontânea a uma porção; o número de porções por disparo segue a lei de Poisson. & \cite{delCastillo1954} & medido \\
3 & O receptor do destinatário é um detector de coincidência; o cálcio desencadeia a mudança de peso & Patch-clamp de neurônios de camundongo: íons Mg$^{2+}$ bloqueiam o canal aberto pelo glutamato no potencial de repouso e saem com a despolarização. Fatias de hipocampo: um quelante de cálcio no destinatário bloqueia a potenciação de longo prazo, e a liberação de cálcio por luz dentro da célula por si só reforça a transmissão. & \cite{Nowak1984, Malenka1988calcium} & medido \\
4 & Qualquer lado executa a decisão: cresce $A_1$ no destinatário, muda $p$ no remetente & Glutamato liberado por luz sobre uma única espinha causa um crescimento rápido e duradouro da espinha e da corrente por seus receptores; a potenciação é acompanhada da inserção de receptores AMPA. A despolarização do destinatário suprime temporariamente a liberação no remetente; o efeito é levado por endocanabinoides que voltam pela fenda (mostrado em entradas \nt{GABA}). A depressão de curto prazo depende da probabilidade de liberação no remetente. & \cite{Matsuzaki2004, Malinow2002, WilsonNicoll2001, Tsodyks1997} & medido \\
5 & A sinapse se reforça quando remetente e destinatário estão ativos juntos~\eqref{eq:hebb} & Coelho anestesiado: após séries curtas de disparos na via de entrada do hipocampo, a resposta fica reforçada por $30$~min a $10$~h. Em fatia de hipocampo, os disparos do destinatário reforçam só as sinapses ativas no mesmo momento. & \cite{Bliss1973LTP, Kelso1986hebb} & medido \\
6 & A regra~\eqref{eq:oja} encontra o autovetor principal das correlações das entradas (proposição~\ref{prop:oja}) & Teorema; demonstração no capítulo. Não há experimento em que um neurônio tenha aprendido a componente principal de suas entradas; o mecanismo que limita o crescimento dos pesos na sinapse não foi estabelecido. & \cite{Oja1982} & teoria \\
7 & Hebb temporal: janela de cerca de $\pm20\ms$ (fig.~\ref{fig:stdp}); sequências repetidas fazem crescer cadeias & Pares de neurônios de hipocampo em cultura: um disparo do destinatário até $20\ms$ depois do disparo do remetente produz reforço duradouro, até $20\ms$ antes, enfraquecimento; ambos dependem de receptores NMDA. A razão $A_-=0.6A_+$ na figura é ilustrativa. Que assim crescem cadeias na rede não foi medido diretamente. & \cite{BiPoo1998} & medido \\
8 & Traço~\eqref{eq:threefactor}: a dopamina que chega $1$--$2\,\text{s}$ após a atividade conjunta reforça a conexão; mais tarde, não (proposição~\ref{prop:threefactor}) & Fatias de estriado, estimulação optogenética separada das entradas glutamatérgicas e dopaminérgicas: a espinha só cresce se a dopamina chega $0.3$--$2\,\text{s}$ após a entrada glutamatérgica; a janela é dada pela subida e queda rápidas do AMPc. No córtex, o pareamento deixa traços separados para reforço e enfraquecimento, que receptores de monoaminas convertem em mudanças duradouras numa janela da ordem de segundos. & \cite{Yagishita2014, He2015eligibility} & medido \\
9 & Janelas de segundos também existem sem dopamina: o terceiro sinal é uma forte excitação do destinatário & Camundongo vivo, campo CA1: um único evento de platô no destinatário reforça as entradas chegadas segundos antes e depois dele e cria um campo de lugar numa só passagem. & \cite{Bittner2017} & medido \\
10 & O passo médio da regra de três fatores segue o gradiente da recompensa~\eqref{eq:perturbation} (proposição~\ref{prop:gradient}) & Teorema do gradiente para aprendizado por desvios aleatórios; deduzido no capítulo para ruído pequeno. & \cite{Williams1992} & teoria \\
11 & O ruído é busca: a rede aprende com seus próprios desvios aleatórios & Mandarins jovens: desligar o núcleo de saída do circuito dos gânglios da base elimina, de forma abrupta e reversível, a variabilidade do canto. Mandarins adultos: uma interferência é aplicada às execuções de uma sílaba com altura de um lado da média, e as aves deslocam rapidamente a altura para o outro lado, aprendendo com a própria dispersão. & \cite{Olveczky2005, TumerBrainard2007} & medido \\
12 & A escolha entre duas opções é aprendida com $\tau_e=1\,\text{s}$ e $\delta=R-\bar R$; não com $\tau_e=50\ms$; sem a subtração, os pesos crescem sem limite, e com taxa de aprendizado alta a rede pode fixar a pior escolha & \texttt{models/dofa\_learning.py}, $\eta=0.05$, $500$ tentativas, $6$ sementes. $\tau_e=1\,\text{s}$: fração de escolhas de $A$ $0.65$--$0.79\to0.96$--$1.00$, pesos $0.85$--$1.33$. $\tau_e=50\ms$: $0.38$--$0.55$, os pesos não mudam. $\delta=R$: peso do vencedor $860$--$1190$. $\delta=R$, $\eta=0.5$, $3$ sementes: uma se fixou em $B$ ($0.04\to0$). O script imprime os quatro casos; o recálculo coincidiu com o capítulo. & \texttt{models/\allowbreak dofa\_\allowbreak learning.py} & modelo \\
13 & O tempo de aprendizado por um único sinal comum cresce proporcionalmente ao número de neurônios ruidosos (proposição~\ref{prop:broadcastcost}) & Curvas de aprendizado de uma rede linear: a perturbação de neurônios é mais lenta que o gradiente exato tantas vezes quantos neurônios de saída houver; a perturbação de pesos, ainda tantas vezes quantas entradas houver. & \cite{Werfel2005} & teoria \\
14 & As saídas do \nt{DOPA} vão sobretudo às regiões de escolha de ações & Reconstrução completa dos axônios de neurônios \nt{DOPA} nigroestriatais isolados do rato: os ramos de uma célula cobrem $0.45$--$5.7\,\%$ do volume do estriado (em média $2.7\,\%$); fora do estriado quase não há ramos. & \cite{Matsuda2009} & medido \\
15 & O córtex aprende a prever sua entrada, e o erro é calculado no local & Revisão: no córtex sensorial há respostas ao desacordo entre a entrada esperada e a recebida. Que isso seja uma regra geral de aprendizado do córtex não foi demonstrado. & \cite{Keller2018} & hipótese \\
16 & A dopamina se comporta como o erro~\eqref{eq:ctd}: pico, transferência para o sinal antecipador, queda (fig.~\ref{fig:ctdcases}) & Neurônios \nt{DOPA} de macaco: pico com suco inesperado; após o aprendizado, pico no sinal antecipador e nenhuma resposta ao suco prometido; queda quando o suco não veio. A forma contínua~\eqref{eq:ctd} é teoria. & \cite{Schultz1997, Doya2000} & medido \\
17 & Circuito que calcula $\dd V/\dd t$ e subtrai a expectativa & Camundongos, registro e optogenética na área tegmental ventral: os neurônios \nt{DOPA} subtraem a expectativa da recompensa; neurônios \nt{GABA} vizinhos os inibem quando a recompensa é esperada, e excitá-los ou inibi-los muda o erro. Como a derivada é calculada não foi mostrado. & \cite{Eshel2015arithmetic} & hipótese \\
18 & O neurônio \nt{DOPA} é um marca-passo; as quedas são menores que os picos & Em fatia, os neurônios \nt{DOPA} descarregam espontaneamente e de modo muito regular, sobre um potencial marca-passo lento. No macaco, a taxa acompanha quantitativamente o erro positivo e transmite o negativo apenas fracamente. & \cite{GraceOnn1989, Bayer2005} & medido \\
19 & Ator e crítico (fig.~\ref{fig:actorcritic}) & O algoritmo vem da teoria do aprendizado por reforço. Em humanos (fMRI), o estriado ventral se comporta como crítico com e sem escolha; o dorsal, só quando é preciso escolher ações. & \cite{SuttonBarto2018, ODoherty2004actor} & hipótese \\
20 & O sinal de $\delta$ na regra é invertido nos neurônios com a segunda família de receptores & Fatias de estriado: nos neurônios com receptores D1, o pareamento produz reforço, que requer D1; nos neurônios com D2, enfraquecimento, que requer D2. & \cite{Shen2008} & medido \\
\bottomrule
\end{longtable}}
